A message of the stream is a record:
the order $(t, q, p, d)$, and what is to be done with it.
A class of records needs a constructor that stores its arguments,
a printed form that shows them,
and an equality that compares them,
and the three are the same for every such class, up to the names of the fields.
A \emph{decorator} is written above a definition, after the sign \codehl{@}:
it is a function that receives the class or the function defined below it,
and returns the object that the name is then bound to.
The decorator \codehl{dataclass} writes the three from the list of the fields
\citep[Chapter 5]{Ram22flu},
as the special methods \codehl{\_\_init\_\_}, \codehl{\_\_repr\_\_} and \codehl{\_\_eq\_\_}
of Table \ref{tab.operationsAsCalls}.
Listing \ref{listing.messageRecord} is the class \codehl{Message}.

\begin{lstlisting}[caption={\protect\codehl{unito26.lob.messages.Message}, abridged: the docstring and two of its three checks are cut.}, label=listing.messageRecord]
@dataclass(frozen=True, slots=True)
class Message:
    # ...
    time: float
    size: int
    price: int
    direction: int
    kind: MessageType = MessageType.SUBMIT

    def __post_init__(self) -> None:
        if self.direction not in (BUY, SELL):
            # ...
        if self.size < 0:
            raise ValueError(f"size must be non-negative, got {self.size}")
        # ...
\end{lstlisting}

The two arguments of the decorator change what an instance is.
With \codehl{frozen=True} an assignment to a field raises an error,
so a message is immutable in the sense of Definition \ref{def.object}:
the stream is folded more than once,
and a message edited in passing would make the second fold differ from the first.
Two messages with equal fields are then equal, though not identical;
and since a frozen data class also receives \codehl{\_\_hash\_\_},
they are one element of a set.
With \codehl{slots=True} the attributes of an instance are its fields and no others,
and Remark \ref{remark.misspeltAttribute} does not apply.

The generated constructor ends by calling \codehl{\_\_post\_init\_\_}, when the class binds it.
That of \codehl{Message} raises on a direction other than $\pm 1$,
on a negative size and on a negative price,
so no instance of \codehl{Message} carries any of the three.
Listing \ref{listing.messageValue} checks these statements;
in it, \codehl{with pytest.raises(E)} asserts that the statement under it raises the error \codehl{E}.
It also checks that two books in the same configuration are not equal:
\codehl{AggregateBook} binds no \codehl{\_\_eq\_\_},
and \codehl{object} compares by identity.

\begin{lstlisting}[caption={A frozen data class is a value, and a plain class is not.}, label=listing.messageValue]
from dataclasses import FrozenInstanceError

import pytest

from unito26.lob.messages import BUY, Message
from unito26.lob.orderbook import AggregateBook

message, twin = Message(1.0, 100, 99, BUY), Message(1.0, 100, 99, BUY)
assert message == twin and message is not twin
assert len({message, twin}) == 1
with pytest.raises(FrozenInstanceError):
    message.size = 0
with pytest.raises(ValueError):
    Message(1.0, 100, 99, 0)

book = AggregateBook.from_levels({99: 100}, {101: 100})
other = AggregateBook.from_levels({99: 100}, {101: 100})
assert book != other and book.levels_map(BUY) == other.levels_map(BUY)
\end{lstlisting}

\begin{remark}\label{remark.frozenIsShallow}
 A frozen instance cannot rebind its fields,
 and the objects they refer to keep their own mutability,
 as Definition \ref{def.object} allows.
 \codehl{SubmitResult} is frozen, and the list of fills it holds can be appended to.
 \codehl{HawkesParams} is frozen,
 and its constructor refuses a branching ratio $\branchingRatio \geq 1$;
 its arrays can still be written to once it has returned.
\end{remark}

\begin{remark}\label{remark.equalityOfArrays}
 The generated \codehl{\_\_eq\_\_} compares the fields with \codehl{==}.
 Between two arrays \codehl{==} returns an array, one entry for each pair of entries,
 and not a truth value.
 So the tests of the package compare two instances of \codehl{HawkesParams}
 field by field, with \codehl{numpy.allclose},
 and \codehl{MarketSession}, whose fields are tables, is declared with \codehl{eq=False}.
\end{remark}
